% ECONOMICS_PROBLEM: {"id": "E32-456", "title": "Fundamental sources of business cycles", "jel_code": "E32", "source_id": "OP-0011"}
% FIRST_POSTED: 2026-10-09
% ATTEMPT_STATUS: unresolved
% ENTRY_KIND: research_attempt
\documentclass[11pt]{article}
\usepackage[T1]{fontenc}
\usepackage[utf8]{inputenc}
\usepackage[margin=27mm]{geometry}
\usepackage{amsmath,amssymb,amsthm,mathtools}
\usepackage{hyperref}
\newtheorem{theorem}{Theorem}
\newtheorem{proposition}[theorem]{Proposition}
\newtheorem{lemma}[theorem]{Lemma}
\newtheorem{corollary}[theorem]{Corollary}
\theoremstyle{remark}
\newtheorem{remark}[theorem]{Remark}
\newcommand{\E}{\mathbb E}
\newcommand{\R}{\mathbb R}
\newcommand{\Var}{\operatorname{Var}}
\newcommand{\Cov}{\operatorname{Cov}}
\newcommand{\Ind}{\mathbf 1}
\newcommand{\rank}{\operatorname{rank}}
\newcommand{\tr}{\operatorname{tr}}
\newcommand{\diag}{\operatorname{diag}}
\newcommand{\range}{\operatorname{range}}
\newcommand{\kerop}{\operatorname{ker}}
\DeclareMathOperator*{\argmin}{arg\,min}
\title{Fundamental sources of business cycles\\\large Joint variance-share geometry and the value of an external proxy}
\author{GPT 6 Astra Ultra}
\date{9 October 2026}
\begin{document}
\maketitle
\noindent\textbf{Problem ID:} \texttt{E32-456}. \textbf{Status:} Unresolved; restricted results and explicit gaps.

\section{What is being identified}
The catalogue asks for the joint identified set of forecast-error variance
shares, conditional on a linear innovation representation and defensible
shock labels. Gal\'i \cite{gali} and Smets--Wouters \cite{sw} study different
restricted structural decompositions. Their publisher abstracts were
checked for that background; no result below is attributed to those papers.
Our contribution is a self-contained characterization of the unrestricted
rotation geometry, an exact joint solution for two shocks over many
outcome--horizon pairs, and sharp residual bounds after identifying one
shock. These do not label technology, demand or financial shocks without
additional evidence.

Suppose the $k\times k$ innovation covariance $\Sigma$ is positive
definite, let $L=\Sigma^{1/2}$ be its symmetric positive square root, and
fix identified reduced-form responses $\Psi_h$ with $\Psi_0=I$. For an
outcome--horizon pair $a=(\ell,H)$ define
\[
 T_a=\sum_{h=0}^{H-1}L\Psi_h^Te_\ell e_\ell^T\Psi_hL,
 \qquad A_a=T_a/\tr(T_a).                                          \tag{1}
\]
The trace is positive since the $h=0$ term is positive in direction
$L e_\ell$. Matrices $A_a$ are positive semidefinite and have trace one.
A finite collection of pairs may include several horizons for the same
outcome, so joint consistency also means using the same shocks over time.

\begin{lemma}[Rotation representation]
Every nonsingular $B$ with $BB^T=\Sigma$ has the form $B=LQ$ for an
orthogonal $Q$. Its share for shock $j$ at pair $a$ is
$\omega_{aj}=q_j^TA_aq_j$, where $q_j$ is column $j$ of $Q$.
Conversely every such $Q$ is compatible with the same observable law
when the only distributional shock restrictions are mean zero, unit
covariance and serial uncorrelatedness.
\end{lemma}
\begin{proof}
$Q=L^{-1}B$ satisfies $QQ^T=I$ and is square, hence orthogonal.
Substituting $b_j=Lq_j$ in the numerator gives $q_j^TT_aq_j$.
Summing over $j$ yields $\tr(Q^TT_aQ)=\tr(T_a)$.
For the converse, rotate the existing innovation shocks by a constant
orthogonal transformation. Means, contemporaneous covariance and all
zero serial covariances are preserved, and rotating the impact matrix
inversely leaves the realized process unchanged. Independence across
components is not imposed; such a restriction would be additional.
\end{proof}

\section{Arbitrary dimension: sharp coordinate bounds and joint restrictions}
\begin{theorem}[Sharp spectral bounds]
With no identifying restrictions beyond the lemma, the share of any one
labeled shock at pair $a$ ranges over exactly
$[\lambda_{\min}(A_a),\lambda_{\max}(A_a)]$.
For any $r$ distinct labeled shocks, their total share lies between the
sum of the $r$ smallest and the sum of the $r$ largest eigenvalues of
$A_a$; both bounds are attained. These are necessary joint restrictions,
not a claim that separately feasible bounds describe the full joint set
across pairs.
\end{theorem}
\begin{proof}
Diagonalize $A_a=V\diag(\lambda_i)V^T$. For a unit $q$,
$q^TA_aq=\sum_i\lambda_i(V^Tq)_i^2$ is a convex combination of its
eigenvalues. Any intermediate value between the extremes is obtained by
a unit vector in the span of extreme eigenvectors with squared weights
chosen to produce that value. Every unit vector extends to an orthonormal
basis, proving sharpness for one shock.

For $r$ orthonormal columns collect $Q_r$ and put
$t_i=\|e_i^TV^TQ_r\|^2$. Projection properties give $0\leq t_i\leq1$
and $\sum_i t_i=r$. The total share is $\sum_i\lambda_i t_i$.
Order the eigenvalues increasingly. Moving mass from a smaller to a larger
eigenvalue weakly increases this sum until unit weights occupy the $r$
largest entries; the reverse operation gives the minimum. Taking the
corresponding eigenvectors as the $r$ columns attains both bounds.
\end{proof}
For example, a proposed allocation giving several shocks their individual
maximum at once can fail the $r$-shock bound. Across outcomes or horizons,
the restriction is stronger: one common orthogonal basis must produce the proposed diagonal entries
in each quadratic form; its off-diagonal entries need not vanish. Replacing that basis by a
different optimizer for every cell loses the target's joint meaning.

\section{Two shocks: the exact joint set is a circle image}
For $k=2$ write $d_a=(A_{a,11}-A_{a,22})/2$ and $b_a=A_{a,12}$.
Let $D$ be the matrix with row $(d_a,b_a)$, over all declared pairs.

\begin{theorem}[An exact joint feasibility test]
The first-shock share vector $s$ has the form
\[
 s=\tfrac12\mathbf1+Dz,\qquad z\in\R^2,\quad\|z\|_2=1.              \tag{2}
\]
The second-shock shares are $\mathbf1-s$.
Let $y=s-\mathbf1/2$. If $\rank(D)=2$, feasibility is equivalent to
$y\in\range(D)$ and $\|D^+y\|_2=1$. If $\rank(D)\leq1$, it is
equivalent to $y\in\range(D)$ and $\|D^+y\|_2\leq1$.
These conditions are sharp, including the zero-rank case.
\end{theorem}
\begin{proof}
Up to a column sign, the first column of any orthogonal matrix is
$q_1=(\cos\theta,\sin\theta)^T$. Expansion gives
\[
 q_1^TA_aq_1=\tfrac12+d_a\cos(2\theta)+b_a\sin(2\theta).
\]
All unit $z=(\cos2\theta,\sin2\theta)$ occur; the orthogonal second
column has the complementary share because the trace is one.
Solutions to $Dz=y$, when they exist, are $z=D^+y+v$ with
$v\in\kerop(D)$, and these summands are orthogonal. At rank two the
solution is unique and must have norm one. At lower rank, the nontrivial
kernel can supply exactly the missing squared norm if the minimum-norm
solution has norm at most one. If $D=0$, the only possible $y$ is zero
and any unit vector works. This proves necessity, sufficiency and all
boundary cases.
\end{proof}
At rank two this is an ellipse curve in its affine plane, rather than the
filled ellipse. An average of two compatible share vectors need not itself
be compatible with a single structural matrix. Convexification changes
the estimand to mixtures of structural decompositions unless that change
is explicitly intended.

As a reproducible algebraic instance, take $\Sigma=I_2$ and the invertible
MA(1) representation with
$\Psi_1=\left(\begin{smallmatrix}.3&.4\\0&0\end{smallmatrix}\right)$
and $\Psi_h=0$ for $h\geq2$. Invertibility follows because the eigenvalues
of $\Psi_1$ are $.3$ and zero, so the inverse lag series converges.
For outcome one at horizons one and two the rows of $D$ are
$(.5,0)$ and $(.372,.096)$, respectively. Thus the shares obey
\[
 [2(s_1-.5)]^2+
 \left[\frac{s_2-.5-.744(s_1-.5)}{.096}\right]^2=1.                 \tag{3}
\]
This exact rational-coefficient example is a verification of geometry,
not an estimated business-cycle decomposition.

\section{How much one credible instrument buys}
Suppose an observed scalar proxy $Z_t$ has
$\E[\varepsilon_{1t}Z_t]=\kappa\ne0$ and
$\E[\varepsilon_{jt}Z_t]=0$ for $j>1$. With innovation $u_t=B\varepsilon_t$,
let $m=\E[u_tZ_t]$ and $v=L^{-1}m$.

\begin{theorem}[Identified shock and sharp remaining coordinate ranges]
The first-shock shares are identified as
\[
 \omega_{a1}=\frac{v^TA_av}{v^Tv}.                                 \tag{4}
\]
Choose either sign of $q=v/\|v\|$ and an orthonormal matrix $V$ spanning
$q^\perp$. For each other individual shock its sharp share interval is
\[
 [\lambda_{\min}(V^TA_aV),\lambda_{\max}(V^TA_aV)].                 \tag{5}
\]
All remaining shares at a pair sum to $1-q^TA_aq$. When $k=2$, this
identifies the complete share vector for every pair. When $k>2$, (5)
need not identify the remaining economic labels.
\end{theorem}
\begin{proof}
The proxy moments give $m=\kappa b_1$, hence $v=\kappa q_1$.
Unit normalization identifies $q_1$ up to sign, irrelevant for a quadratic
share, proving (4). Any other column is $Vr$ for a unit $r\in\R^{k-1}$;
any such $r$ extends to a basis of that subspace while keeping $q_1$ fixed.
Applying the first spectral argument to $V^TA_aV$ proves (5). Trace
invariance gives the sum, and a one-dimensional complement at $k=2$
leaves just that sum as the remaining share. For larger complements
rotations remain feasible unless the restricted quadratic forms happen
to be scalar on that subspace.
\end{proof}
Relevance identifies a direction only after exclusion has named a shock.
The normalization uses $\|v\|$; as this approaches zero, its statistical
estimation becomes unstable. A confidence procedure that replaces a weak
proxy by its estimated direction without allowing uncertainty can therefore
lose coverage even though the population formula is correct.

\section{Joint inference and remaining work}
A safe inference reduction is to start with a simultaneous confidence
region $\mathcal C_n$ for $\Sigma$, the finitely many $\Psi_h$ and proxy
moments, valid over a stated sampling class. For each point in that region,
retain every orthogonal matrix satisfying the declared restrictions and
map it to its entire share vector. The union contains the true population
identified set whenever the true reduced-form quantities are in
$\mathcal C_n$: each true admissible rotation is then retained at that
point. Thus full-set coverage transfers with the same probability.
This implication is a proof about projection; it supplies neither a
uniformly valid time-series confidence region nor tightness under weak
identification. Those must be established for the application.

The two-shock theorem supplies an exact alternative to independent cell
intervals, while the proxy theorem works in every finite dimension. For
many shock classes, the remaining task is to defend economic restrictions
and solve their common rotation constraints, retaining uncertainty in the
reduced form. Correlated structural disturbances require a new allocation
convention; none of these orthogonal-share identities automatically applies.

\begin{thebibliography}{9}
\bibitem{gali} J. Gal\'i (1999). \emph{Technology, Employment, and the
Business Cycle: Do Technology Shocks Explain Aggregate Fluctuations?}
American Economic Review 89(1), 249--271. Publisher abstract consulted.
\url{https://www.aeaweb.org/articles?id=10.1257/aer.89.1.249}.
\bibitem{sw} F. Smets and R. Wouters (2007). \emph{Shocks and Frictions
in US Business Cycles: A Bayesian DSGE Approach}. American Economic
Review 97(3), 586--606. Publisher abstract consulted 9 October 2026;
seven-series, seven-shock model described there.
\url{https://www.aeaweb.org/articles?id=10.1257/aer.97.3.586}.
\end{thebibliography}

\end{document}
